%% SCRAP: papers/ANTI_CLAIMS
%% SOURCE: docs/working/papers/ANTI_CLAIMS.md
%% STATUS: CURRENT
%% FITS: ssrn/ch-anti-claims, vol3-research/ch-formal-claims
%% EDITORIAL: lifted — prose rewritten to press voice

\section{Out-of-Scope Claims: What This Work Does Not Assert}

This section enumerates claims the StarForth project explicitly does \emph{not} make.
Its purpose is to bound the scope of peer review: a reviewer who attributes an
unstated claim to the work is engaging a strawman. Each entry pairs what
the implementation \emph{does} assert with what it \emph{does not} assert.
The final summary table (\S\ref{sec:anti-claims-summary}) is the canonical
reference for authors responding to review challenges.

\subsection{Physics and Thermodynamics}

\paragraph{Not claimed: a physical theory.}
The adaptive runtime uses execution frequency as a proxy for thermal energy
and exponential decay as a model of heat dissipation. These are conceptual
tools borrowed from thermodynamics, not claims that physical laws govern code
execution. No actual thermal processes occur in the CPU as a consequence of
this model; no quantum effects are invoked; energy conservation in the
thermodynamic sense does not apply to execution frequency.

Precise language: \emph{``execution frequency evolves like heat in a cooling
system,''} not \emph{``execution frequency is heat.''}

\paragraph{Not claimed: frequency is thermal energy.}
The implementation uses a 64-bit integer counter (\texttt{uint64\_t
execution\_heat}) incremented on each word execution. It does not measure
joules, calories, or CPU die temperature. The term ``heat'' is a metaphorical
label. Exponential decay \emph{resembles} heat dissipation and
steady-state convergence \emph{mirrors} thermodynamic equilibrium in a
structural, mathematical sense only.

\subsection{Optimality and Performance}

\paragraph{Not claimed: global optimality.}
The system demonstrates a 25.4\% performance improvement over baseline under
the conditions described. This does not imply that no other adaptive runtime
could perform better, that optimality is proven mathematically, or that the
implementation outperforms all other virtual machines. The contribution is
\emph{a working adaptive system}, not the globally optimal one.

\paragraph{Not claimed: superiority to JIT compilers.}
JIT compilers such as PyPy, LuaJIT, and HotSpot share the conceptual goal of
frequency-based specialization. The claim here is not superior raw speed but
\emph{deterministic adaptation}: the same adaptive decisions occur on every
run, enabling formal verification of the optimization mechanism itself. JITs
typically cannot make this guarantee.

\paragraph{Not claimed: universal workload applicability.}
The system is validated on CPU-bound, deterministic FORTH programs, including
recursive algorithms (Fibonacci, Ackermann) across workload shapes with Zipf
exponent $\alpha \in [0.8, 1.5]$. It does not claim to improve I/O-bound
workloads, non-deterministic programs, adversarial execution patterns, or very
short-lived processes (fewer than approximately 1{,}000 iterations). Explicit
failure modes are documented in the companion negative-results section.

\subsection{Machine Learning and Artificial Intelligence}

\paragraph{Not claimed: AI or machine learning.}
The adaptive mechanism uses statistical inference (ANOVA, Levene's test,
exponential regression) and data-driven parameter tuning. No neural networks,
gradient descent, backpropagation, training datasets, deep learning, or
reinforcement learning are involved.

Precise language: \emph{``statistically-inferred adaptive tuning,''} not
\emph{``AI-driven optimization.''}

\paragraph{Not claimed: the system learns.}
The system \emph{adapts}: parameters converge to a steady state and feedback
loops stabilize metrics. This is statistical convergence, not supervised or
unsupervised learning. Knowledge does not transfer between workloads; the
steady state is workload-specific.

\subsection{Novelty and Prior Art}

\paragraph{Not claimed: first adaptive virtual machine.}
Execution frequency tracking is standard profiler practice; hot-code caching
is the basis of every production JIT; exponential decay underlies LRU
eviction. The novelty claim is specific: the \emph{combination} of adaptation
with 0\% algorithmic variance and formal verification potential, which prior
systems do not offer.

\paragraph{Not claimed: a replacement for JIT compilation.}
The implementation targets a different niche---verifiable adaptive systems
for safety-critical contexts---rather than competing directly with LLVM or
V8. Compilation remains necessary for applications requiring peak throughput.

\subsection{Formal Verification}

\paragraph{Not claimed: complete formal verification.}
% PATENT: formal verification claims intersect patent scope; do not strengthen
Convergence theorems are stated and empirically validated (0\% coefficient of
variation across 90 runs). Full mechanized proofs for all seven feedback loops
in Coq or Isabelle are \emph{not} claimed; partial proofs are in progress.

Precise language: \emph{``empirically validated determinism,''} not
\emph{``formally proved correct.''}

\paragraph{Not claimed: zero bugs.}
The test suite comprises 936\raisebox{0.5ex}{+} tests and validates FORTH-79
compliance. No known correctness bugs exist in the core interpreter as of the
experimental baseline commit. This provides high assurance, not mathematical
proof of defect-absence.

\subsection{Causation and Interpretation}

\paragraph{Not claimed: proven causation.}
The association between adaptive mechanisms and the observed 25.4\%
improvement is strong: the functional relationship fits an exponential decay
model with $R^2 > 0.95$, and only the fully adaptive configuration
(\texttt{C\_FULL}) improves. No randomized controlled trial establishes
mechanistic causation.

Precise language: \emph{``adaptive mechanisms are associated with 25.4\%
improvement,''} not \emph{``cause 25.4\% improvement.''}

\paragraph{Not claimed: a theory of computation.}
The work provides a VM optimization technique and a predictive framework for
adaptive runtime performance. It does not propose a fundamental theory of
computation, replace computational complexity theory, or define a new model
of computation.

\subsection{Generalization}

\paragraph{Not claimed: cross-language generalization.}
Principles may generalize to other interpreter architectures (Lua, Python,
JavaScript), but this has not been validated. Compiled languages are
explicitly out of scope: ahead-of-time optimization already handles hot-code
without runtime frequency tracking.

\paragraph{Not claimed: arbitrary scalability.}
The system is validated on programs up to approximately 2.1 million word
executions with dictionary sizes up to approximately 500 entries. Scalability
to programs with 100{,}000\raisebox{0.5ex}{+} dictionary entries is not
tested; transition-matrix memory overhead would grow quadratically in that
regime.

\subsection{Deployment Status}

\paragraph{Not claimed: production readiness.}
StarForth is a research prototype suitable for experimental validation. It
demonstrates feasibility of deterministic adaptation and serves as a platform
for further study. It is not a production-grade implementation with enterprise
support or hardening against all security threats.

\paragraph{Not claimed: an operating system.}
The StarForth $\to$ StarKernel $\to$ StarshipOS sequence is a roadmap, not a
set of delivered products. StarshipOS does not yet exist as a functional
system; the vision is aspirational.

\subsection{Statistical Claims}

\paragraph{Not claimed: zero variance in all metrics.}
The 0.00\% coefficient of variation applies to \emph{algorithmic decisions}:
cache hit rates and dictionary lookup paths. Wall-clock runtime exhibits
60--70\% CV due to OS scheduler noise, thermal variation, and cache-line
effects. These two components are statistically independent (Pearson
$r = 0.03$, $p = 0.87$).

Precise language: \emph{``algorithmic variance: 0.00\% CV,''} not
\emph{``total system variance: 0.00\% CV.''}

\paragraph{Not claimed: 100\% confidence.}
Results are reported at 95\% confidence intervals. The Bayesian posterior
$P(H_1 \mid \text{data}) \approx 1 - 10^{-30}$ is effectively certain but
not unity. Science deals in probabilities, not absolute certainties;
replication failure, while astronomically unlikely, is not logically
impossible.

\subsection{Scope Exclusions}

\paragraph{Not claimed: solutions to undecidable problems.}
Adaptive convergence to steady state is asserted only for \emph{terminating}
programs with analyzable workloads. No claim is made about decidability of
convergence for arbitrary programs.

\paragraph{Not claimed: quantum or blockchain components.}
The implementation uses classical algorithms on classical hardware. No quantum
superposition, entanglement, distributed ledger, cryptocurrency, smart
contracts, or neuromorphic computing is involved.

\paragraph{Not claimed: a performance record over any named system.}
No direct performance shootout against PyPy, LuaJIT, or HotSpot is
presented. The contribution is deterministic adaptation, not a speed record.

\subsection{Summary Table}
\label{sec:anti-claims-summary}

\begin{table}[h]
\centering
\caption{Anti-claims reference table. Left column: what the work asserts.
Right column: what it explicitly does not assert.}
\label{tab:anti-claims}
\begin{tabular}{lll}
\toprule
\textbf{Category} & \textbf{Asserted} & \textbf{Not Asserted} \\
\midrule
Physics         & Thermodynamic metaphor        & Actual physical theory \\
Performance     & 25.4\% improvement            & Global optimality \\
Novelty         & Deterministic adaptation      & First adaptive system \\
Verification    & Empirical validation          & Complete formal proof \\
ML/AI           & Statistical inference         & Neural networks or learning \\
Causation       & Strong correlation            & Proven causation \\
Generality      & Works for FORTH               & Works for all languages \\
Variance        & Algorithmic: 0\% CV           & Total system: 0\% CV \\
\bottomrule
\end{tabular}
\end{table}

\subsection{How to Use This Section in Peer Review}

When a reviewer attributes a claim not appearing in the formal claim table,
the appropriate response is to cite the relevant paragraph above by section
number. If the attributed claim does not appear in this section either, it may
represent a genuine novel claim; in that case, authors should locate the
supporting evidence in the formal claim table before responding.

The intellectual commitment underlying this section is that explicit
scope-bounding prevents both strawman attacks and over-interpretation by
readers. Transparency about limitations strengthens, rather than weakens, the
credibility of the work.
